\documentclass[11pt,a4paper]{article}

% ---------------------------------------------------------
% PREAMBLE (stable, warning-free, Overleaf-safe)
% ---------------------------------------------------------
\usepackage[utf8]{inputenc}
\usepackage[T1]{fontenc}
\usepackage{lmodern}
\usepackage{amsmath, amssymb, amsfonts}
\usepackage{bm}
\usepackage{geometry}
\usepackage{graphicx}
\usepackage{hyperref}
\usepackage{cite}
\usepackage{authblk}
\usepackage{physics}
\usepackage{mathtools}

\geometry{margin=2.4cm}

\title{\textbf{Companion LXXVIII: Schrödinger Bridges and Diffusion-Regularized Transport of Persistence Diagrams in the Relaxation-Driven Cyclic Cosmology}}
\author{Michael Lehmann \\ \texttt{mi.lehmann@gmx.de}}
\date{April 2026}

\begin{document}
\maketitle

\begin{abstract}
Schrödinger bridges provide a dynamic and stochastic generalization of entropic 
optimal transport, describing the most likely diffusion process connecting two 
probability distributions. Applied to persistence diagrams, Schrödinger bridges 
yield diffusion-regularized transport paths that interpolate between topological 
ensembles in a smooth and probabilistically meaningful way. Building on the 
Sinkhorn geometry developed in Companion LXXVII, this work introduces Schrödinger 
bridge dynamics for persistence diagrams in the Relaxation-Driven Cyclic Cosmology 
(RDCC). The relaxation mechanism modifies the temporal evolution of the 
gravitational potential, inducing characteristic changes in diffusion flows, 
stochastic transport costs, and dynamic barycenters. The resulting framework 
provides a powerful tool for analyzing RDCC signatures in weak-lensing topology 
through the lens of stochastic optimal transport.
\end{abstract}
% ---------------------------------------------------------
\section{Introduction}

Persistence diagrams encode the birth and death of topological features in weak-
lensing convergence fields. Classical Wasserstein geometry measures the optimal 
transport cost between diagrams, while entropic optimal transport introduces a 
regularization that yields smooth and computationally efficient Sinkhorn 
divergences.

Schrödinger bridges extend this framework by introducing a dynamic, stochastic 
interpretation: instead of computing a static transport plan, one seeks the most 
likely diffusion process that evolves one distribution of persistence diagrams 
into another. This dynamic viewpoint is particularly relevant for cosmology, where 
the underlying gravitational potential evolves over time.

In the Relaxation-Driven Cyclic Cosmology (RDCC), the relaxation mechanism modifies 
the temporal structure of the gravitational field, producing characteristic 
signatures in diffusion-regularized transport of topological features.
% ---------------------------------------------------------
\section{Schrödinger Bridges}

Given two probability distributions $\rho_{0}$ and $\rho_{1}$ on persistence 
diagrams, the Schrödinger bridge problem seeks the most likely diffusion process 
$(\rho_{t})_{t \in [0,1]}$ connecting them.

The bridge solves

\begin{equation}
    \min_{\rho_{t}, v_{t}}
    \int_{0}^{1}
    \int
        \frac{1}{2} \| v_{t}(D) \|^{2} \rho_{t}(D)
    \, dD \, dt
    + \varepsilon \, \mathrm{KL}(\rho_{t} \| \rho_{\mathrm{ref}}),
\end{equation}

subject to the continuity equation

\begin{equation}
    \partial_{t} \rho_{t}
    + \nabla \cdot (\rho_{t} v_{t}) = 0,
\end{equation}

and boundary conditions



\[
    \rho_{0}, \quad \rho_{1}.
\]



Here $\varepsilon$ controls the diffusion strength.
% ---------------------------------------------------------
\section{Diffusion-Regularized Transport of Persistence Diagrams}

For persistence diagrams, the diffusion process acts on birth-death coordinates:

\begin{equation}
    dX_{t}
    = v_{t}(X_{t}) \, dt
    + \sqrt{2 \varepsilon} \, dB_{t},
\end{equation}

where $B_{t}$ is Brownian motion on the diagram space.

The resulting transport path $(\rho_{t})$ is smoother and more stable than 
classical Wasserstein geodesics, making it well-suited for noisy weak-lensing 
data.
% ---------------------------------------------------------
\section{RDCC Modification of Schrödinger Dynamics}

In RDCC, the convergence evolves as

\begin{equation}
    \kappa_{\mathrm{RDCC}}(\ell)
    = \kappa(\ell)
      \left[
        1 - \chi_{\kappa}
        \left(\frac{\ell}{\ell_{\mathrm{rel}}}\right)^{\alpha}
      \right].
\end{equation}

This modifies the drift term in the Schrödinger bridge:

\begin{equation}
    v^{\mathrm{RDCC}}_{t}(D)
    = v_{t}(D)
      \left[
        1 - \chi_{v}
        \left(\frac{\ell}{\ell_{\mathrm{rel}}}\right)^{\alpha}
      \right].
\end{equation}

The diffusion-regularized transport cost becomes

\begin{equation}
    \mathcal{S}^{\mathrm{RDCC}}
    = \mathcal{S}
      \left[
        1 - \chi_{\mathcal{S}}
        \left(\frac{\ell}{\ell_{\mathrm{rel}}}\right)^{\alpha}
      \right].
\end{equation}
% ---------------------------------------------------------
\section{Dynamic RDCC Signatures}

RDCC induces:

\begin{itemize}
    \item slower diffusion at small scales,
    \item enhanced drift at large scales,
    \item a characteristic turnover near $\ell_{\mathrm{rel}}$,
    \item modified Schrödinger barycenters,
    \item altered stochastic geodesics in topological space,
    \item scale-dependent deformation of diffusion paths across redshift.
\end{itemize}

These signatures provide a dynamic and probabilistic probe of the RDCC 
gravitational field.
% ---------------------------------------------------------
\section{Conclusion}

Schrödinger bridges provide a dynamic and stochastic generalization of entropic 
optimal transport, yielding diffusion-regularized transport paths between 
persistence diagrams. In the Relaxation-Driven Cyclic Cosmology, the relaxation 
mechanism modifies these diffusion flows in a scale-dependent manner, producing 
distinctive signatures in weak-lensing topology. The present Companion establishes 
the foundation for future studies of stochastic optimal transport, diffusion 
models, and dynamic geometric analysis in cosmological topology.

\bibliographystyle{apsrev4-2}
\nocite{*}
\bibliography{references}


\end{document}
